\section{Bias bound under informative coding}
\label{sec:methodology}

The glass ceiling (\cref{thm:glass-ceiling}) shows code data alone
cannot identify prevalence, and chart review
(\cref{thm:identifiability-chart}) restores identifiability under two
conditions. This section quantifies what happens when those conditions
fail. The C2 violation, informative coding, drives the code-only
estimator's bias; the case-mix gap drives the calibrated estimator's
residual bias. We bound each.

\subsection{Setup}

Recall the coding mechanism \eqref{eq:trans-coding}: code sensitivity
for a true case is $\mathrm{logit}^{-1}(b_0 + b_1 S)$, increasing in
latent severity $S$ when $b_1 > 0$. The marginal code sensitivity over
the cohort's diseased patients is
\begin{equation}
  \mathrm{sens}_{\text{coh}}
    = \E[\,\mathrm{logit}^{-1}(b_0 + b_1 S) \mid \Y = 1,\ \text{cohort}\,].
  \label{eq:meth-sens}
\end{equation}
When $b_1 = 0$, coding does not depend on severity and C2 holds. When
$b_1 \neq 0$, C2 fails: the code carries information about $S$ beyond
$\Y$. The natural scalar severity of the violation is the
coding-slope magnitude $|b_1|$ itself, and the bias of the code-only
estimator (below) is driven by $|b_1|$. A descriptive correlation
\begin{equation}
  \rho := \mathrm{Corr}\big(S,\ \mathrm{logit}^{-1}(b_0 + b_1 S)
                              \,\big|\, \Y = 1\big)
  \label{eq:meth-rho}
\end{equation}
is sometimes reported, but $\rho$ is \emph{not} monotone in $|b_1|$
and should not be read as the violation strength: at $b_1 = 0$ the
coding probability is constant and $\rho = 0$ degenerately, whereas
as $b_1 \to 0^{+}$ the logistic is near-linear in $S$ and
$\rho \to 1$; as $|b_1|$ grows the logistic saturates toward a step
function, which \emph{reduces} the linear correlation with $S$.
\Cref{tab:informative} shows $\rho$ rising to $0.994$ at $b_1 = 0.5$
and then falling monotonically to $0.743$ at $b_1 = 3.0$ even as the
code-only bias grows, confirming the non-monotonicity. We therefore
index violation severity by $|b_1|$ throughout.

\subsection{Bias of the code-only estimator}

The code-only estimator $\hat\pi_{\text{code}} = \bar{C}$ targets the
code frequency \eqref{eq:code-freq} rather than the prevalence.

\begin{theorem}[Bias bound under informative coding]
\label{thm:bias-informative}
Under the DGP of \cref{sec:translation}, the code-only prevalence
estimator has expectation
\begin{equation}
  \E[\hat\pi_{\text{code}}]
    = \pi\,\mathrm{sens}_{\text{coh}}
      + (1 - \pi)(1 - \mathrm{spec}),
  \label{eq:meth-codebias}
\end{equation}
so its bias is
\begin{equation}
  \mathrm{Bias}(\hat\pi_{\text{code}})
    = -\pi\big(1 - \mathrm{sens}_{\text{coh}}\big)
      + (1 - \pi)(1 - \mathrm{spec}).
  \label{eq:meth-codebias2}
\end{equation}
Let $\mu_{S,1} := \E[S \mid \Y = 1]$ and $\sigma_{S,1} := \mathrm{sd}(S
\mid \Y = 1)$ be the conditional mean and standard deviation of
severity among diseased patients. The cohort sensitivity admits the
bound
\begin{equation}
  \big|\,\mathrm{sens}_{\text{coh}} \,-\, \mathrm{logit}^{-1}(b_0 + b_1\,\mu_{S,1})\,\big|
    \;\le\; \frac{|b_1|\,\sigma_{S,1}}{4},
  \label{eq:meth-bound}
\end{equation}
so the gap between the average-severity coding probability
$\mathrm{logit}^{-1}(b_0 + b_1\,\mu_{S,1})$ and the marginal cohort
sensitivity is controlled by the severity spread $\sigma_{S,1}$ scaled
by the strength of informative coding $|b_1|$. The sensitivity deficit
$1 - \mathrm{sens}_{\text{coh}}$ is bounded by $1 -
\mathrm{logit}^{-1}(b_0 + b_1\,\mu_{S,1}) + |b_1|\,\sigma_{S,1}/4$.
\end{theorem}

\begin{proof}
Equation \eqref{eq:meth-codebias} is the expectation of a Bernoulli
mixture: a code is positive either because the patient is a true case
who is coded (probability $\pi\,\mathrm{sens}_{\text{coh}}$) or a
non-case who is miscoded (probability $(1-\pi)(1-\mathrm{spec})$).
Subtracting $\pi$ from this expectation gives
\eqref{eq:meth-codebias2}. For \eqref{eq:meth-bound}, write the
per-case code probability as $g(s) := \mathrm{logit}^{-1}(b_0 + b_1
s)$. The logistic derivative $g'(s) = b_1\,g(s)(1-g(s))$ satisfies
$|g'(s)| \le |b_1|/4$ for all $s$, since $x(1-x) \le 1/4$ on $[0,1]$.
Re-anchor at the conditional mean $\mu_{S,1}$ (a population quantity
of the cohort DGP, estimable from the chart-reviewed subsample under
the case-mix representativeness of \cref{thm:identifiability-chart}(ii)):
by the mean value theorem, for every $s$ there exists $\xi$ between
$s$ and $\mu_{S,1}$ with
$g(s) - g(\mu_{S,1}) = g'(\xi)\,(s - \mu_{S,1})$, so
\begin{equation*}
  |g(s) - g(\mu_{S,1})|
    \;\le\;
    \frac{|b_1|}{4}\,|s - \mu_{S,1}|
\end{equation*}
holds pointwise. Taking expectations conditional on $\Y = 1$ and
applying Jensen's inequality to the convex function $|\cdot|$ (so
$|\E[X]| \le \E[|X|]$) gives
\begin{equation*}
  |\,\mathrm{sens}_{\text{coh}} - g(\mu_{S,1})\,|
    \;=\; |\E[g(S) \mid \Y = 1] - g(\mu_{S,1})|
    \;\le\; \frac{|b_1|}{4}\,\E[\,|S - \mu_{S,1}| \,\mid \Y = 1\,].
\end{equation*}
By Jensen's inequality applied to the (concave) square root, $\E[|S -
\mu_{S,1}| \mid \Y = 1] \le \sqrt{\E[(S - \mu_{S,1})^2 \mid \Y = 1]} =
\sigma_{S,1}$. Combining gives \eqref{eq:meth-bound}. The anchor
$g(\mu_{S,1}) = \mathrm{logit}^{-1}(b_0 + b_1\,\mu_{S,1})$ is the
coding probability of a patient at average severity, and the bound
states that the marginal sensitivity over diseased patients departs
from this anchor by at most $|b_1|\sigma_{S,1}/4$. The bound is tight
in the C2-holds direction: when $b_1 = 0$, both sides vanish and
$\mathrm{sens}_{\text{coh}} = \mathrm{logit}^{-1}(b_0)$ exactly.
\end{proof}

Two features of \cref{thm:bias-informative} are worth stating
plainly. First, the code-only bias does not vanish as the cohort grows:
it is a bias in the estimand, not sampling error, exactly the situation
\cref{thm:glass-ceiling} predicts. Second, the \emph{sign} of the bias
is not predetermined. \Cref{eq:meth-codebias2} is a contest between
the sensitivity deficit $\pi(1 - \mathrm{sens}_{\text{coh}})$ and the
false-positive inflation $(1-\pi)(1-\mathrm{spec})$. When coding
sensitivity is low, the deficit dominates and code-only prevalence
\emph{understates} the truth: mild and asymptomatic true cases are
systematically uncoded. When informative coding raises the marginal
sensitivity of the true cases that do present, the deficit shrinks
faster than the false-positive inflation, and for a sufficiently
severity-coded condition the bias turns positive, so code-only
prevalence \emph{overstates} the truth. The bias passes through zero
at the level set $\mathrm{sens}_{\text{coh}} = 1 -
(1-\pi)(1-\mathrm{spec})/\pi$, a coincidence rather than a regular
feature. The simulation of \cref{sec:validation} (\cref{tab:informative})
exhibits this sign change directly: at non-informative coding ($b_1 =
0$) the code-only bias is $-0.016$, and it crosses zero and rises to
$+0.032$ as the severity-coding correlation grows. The unifying
statement is that the bias is structural, governed by the coding
mechanism, and is not a sampling artifact; the direction of the bias
is set by the specific coding-mechanism parameters and is not a useful
heuristic on its own. This is the surveillance-bias mechanism of
\citet{haut2011surveillance,carrell2014surveillance} expressed as a
coarsening statement: more intense detection of severe cases makes the
coded population unrepresentative of the true case population, and
the resulting prevalence error can move in either direction.

\subsection{Residual bias of the chart-review-calibrated estimator}

The calibrated estimator of \cref{thm:identifiability-chart} fits the
coding model, including the severity dependence, on the chart-reviewed
subsample and then deconvolves cohort prevalence via
\eqref{eq:deconvolve}. If the subsample case mix matches the cohort,
condition (ii) of \cref{thm:identifiability-chart}, the calibrated
estimator is consistent. In practice chart-reviewed patients are not a
random sample of the cohort: review is often triggered by an encounter,
a code, or a referral, so the reviewed subsample is enriched for
severe, well-documented patients.

\begin{corollary}[Residual bias from case-mix gap]
\label{cor:casemix}
Let $\mathrm{sens}_{\text{rev}}$ be the marginal code sensitivity in
the chart-reviewed subsample and $\mathrm{sens}_{\text{coh}}$ that in
the cohort, and define the case-mix gap $\delta := \mathrm{sens}_{\text{rev}}
- \mathrm{sens}_{\text{coh}}$. To first order in $\delta$, the
chart-review-calibrated prevalence estimator has bias
\begin{equation}
  \mathrm{Bias}(\hat\pi_{\text{cal}})
    \;\approx\;
    \frac{-\,\pi\,\delta}
         {\,\mathrm{sens}_{\text{coh}} + \mathrm{spec} - 1\,}
    \,+\, O(\delta^2),
  \label{eq:meth-residual}
\end{equation}
so the residual bias vanishes when the reviewed subsample is
representative ($\delta = 0$) and grows linearly in the case-mix gap.
\end{corollary}

\begin{proof}[Sketch]
The calibrated estimator substitutes $\mathrm{sens}_{\text{rev}}$ for
the cohort sensitivity in the deconvolution formula
\eqref{eq:deconvolve}. Differentiating \eqref{eq:deconvolve} with
respect to the plugged-in sensitivity at the true parameters gives the
sensitivity coefficient $-\pi/(\mathrm{sens}_{\text{coh}} +
\mathrm{spec} - 1)$; multiplying by the perturbation $\delta$ yields
\eqref{eq:meth-residual}. The argument is a first-order Taylor
expansion of the deconvolution map, the same perturbation calculation
used for the ERCC-endogenous bias bound in
\citet[\S 7]{towell2026scrnacoarsening}.
\end{proof}

\Cref{cor:casemix} is the phenotyping analog of the ERCC-endogenous
gap. In scRNA-seq, spike-in transcripts have a different capture
efficiency from endogenous transcripts, so the spike-in-calibrated
correction inherits a bias proportional to the capture gap
\citep{towell2026scrnacoarsening}. In phenotyping, chart-reviewed
patients have a different coding sensitivity from the cohort, because
review is triggered by documentation-rich encounters, so the
chart-review-calibrated estimator inherits a bias proportional to the
case-mix gap $\delta$. The structural parallel is exact: a singleton
candidate set restores identifiability only to the extent that the
singletons are drawn from the same population as the units being
corrected.

The practical recommendation that follows mirrors the one in
\citet{towell2026scrnacoarsening}. Chart review is a
\emph{principled but bounded} tool. It is principled because
\cref{thm:identifiability-chart} gives a clean recovery procedure;
bounded because the case-mix gap limits what review can recover.
Practitioners should sample chart review to match the cohort case mix
(for example, stratified review across encounter intensity and code
status) rather than reviewing whoever is most convenient, and should
report the case-mix gap $\delta$ as a sensitivity quantity alongside
the prevalence estimate.

\paragraph{An oracle-free reference.} Calibrated EHR prevalence
estimators that fit code error rates on a chart-reviewed sample are
already available \citep{tong2020augmented,beesley2022samba,zhang2019phiap};
the point of this subsection is narrower, namely that the coarsening
deconvolution can be evaluated without the latent cohort labels, so the
remedy survives outside the simulation oracle. When the fitted
code-sensitivity $\mathrm{logit}^{-1}(\hat b_0 + \hat b_1 S)$ is
marginalized to a scalar $\widehat{\mathrm{sens}}$, the reference
distribution of $S$ over true cases must not use the latent label
$\Y$ for the full cohort, which is unavailable at deployment.
Averaging over $\{S_i : \Y_i = 1\}$ in the whole cohort is an oracle.
The deployable replacement marginalizes over \emph{all} observed
cohort severities weighted by an estimated membership probability
$\widehat\Prob(\Y = 1 \mid S)$, itself fit on the adjudicated
subsample (where $\Y$ is observed). In simulation with $\pi = 0.1$,
$b_1 = 1.5$, and severity-biased review, the oracle reference and the
weighted deployable reference give essentially identical prevalence
estimates (bias $+0.0002$ vs $+0.0004$ over $400$ replicates), while
the convenience subsample reference is biased ($-0.0043$); see
\texttt{scripts/oracle\_check.R}. The chart-review remedy therefore
survives in a form that uses only observed severities and the review
labels.
